\section{这门课能教你什么}

讲师把课程知识分成三类：mechanics、mindset 和 intuitions。

\begin{table}[H]
\centering
\begin{tabular}{p{2.6cm}p{4.8cm}p{6.2cm}}
\toprule
\textbf{类别} & \textbf{含义} & \textbf{可迁移性} \\
\midrule
Mechanics & 机制层知识：Transformer 是什么、并行化怎么做、系统如何配合硬件 & 高，可系统训练 \\
Mindset & 规模意识：如何尽可能压榨硬件、如何认真对待 scaling law & 高，可明确传递 \\
Intuitions & 哪些数据和建模决策真的有效 & 部分可传递，强依赖尺度与任务 \\
\bottomrule
\end{tabular}
\end{table}

其中 mechanics 和 mindset 是这门课最稳定、最值钱的部分。
intuitions 当然也重要，但它更依赖经验、实验和尺度环境，未必能从一个尺度直接外推到另一个尺度。

\subsection{经验和实验，有时比解释更诚实}

在字幕里，讲师提到很多 Transformer 里的设计，未必都有漂亮的第一性原理解释。
比如 SwiGLU 这类非线性，很多时候就是“实验结果确实好”，然后被采用。
课程希望传达的不是“我们已经完全解释了所有设计”，而是“实验和规模经常比故事更可靠”。

\begin{importantbox}{一句务实的话}
如果某个设计选择能在更大规模上稳定提升效果，那么它就值得重视，哪怕它一开始并没有完整的理论解释。
\end{importantbox}

